% --- Archon physics formalization source begin ---
% archon:physics
% archon:covers PhyXMiniProblems/problem_phyx_mini_0917.lean
% archon:source-report reports/phyx_mini/problem_phyx_mini_0917.source.json

\chapter{Physics problem phyx\_mini\_0917}
\label{ch:PhyXMiniProblems_problem_phyx_mini_0917}

\paragraph{Problem source.}
\#\# Physical scenario

The magnetic field between the poles of the electromagnet in figure is uniform at any time, but its magnitude is increasing at the rate of \$0.020 \textbackslash{}, \textbackslash{}text\{T/s\}\$. The area of the conducting loop in the field is \$120 \textbackslash{}, \textbackslash{}text\{cm\}\textasciicircum{}2\$, and the total circuit resistance, including the meter, is \$5.0 \textbackslash{}, \textbackslash{}Omega\$.

\#\# Question

Find the induced current in the circuit.

\#\# Answer choices

- A: A: 0.024mA

- B: B: -0.024mA

- C: C: 0.048mA

- D: D: -0.048mA

\#\# Figure caption (auxiliary; use the image as primary evidence)

The image depicts a setup involving electromagnetic induction. Here's a breakdown of the components and elements:

1. **Components**:
   - Two parallel plates labeled "S" (upper) and "N" (lower), representing the poles of a magnet.
   - A circular loop positioned horizontally between the poles labeled "a" and "b." 
   - A connected meter shown measuring current, with a pointer indicating a reading.
   - The magnetic field is represented by vertical blue arrows flowing from "S" to "N."

2. **Labels and Text**:
   - \textbackslash{}( \textbackslash{}frac\{dB\}\{dt\} = 0.020 \textbackslash{}, \textbackslash{}text\{T/s\} \textbackslash{}): Denotes the rate of change of the magnetic field.
   - \textbackslash{}( A = 120 \textbackslash{}, \textbackslash{}text\{cm\}\textasciicircum{}2 = 0.012 \textbackslash{}, \textbackslash{}text\{m\}\textasciicircum{}2 \textbackslash{}): Indicates the area of the loop.
   - "Total resistance in circuit and meter = 5.0 \textbackslash{}( \textbackslash{}Omega \textbackslash{})": Specifies the total resistance in the circuit.

3. **Vectors and Directions**:
   - A vector \textbackslash{}( \textbackslash{}vec\{A\} \textbackslash{}) is shown pointing upward, indicating the direction of the area of the loop.
   - A red arrow labeled "I" points to the left, representing the direction of the induced current in the loop.

4. **Relationships**:
   - The setup illustrates electromagnetic induction where a changing magnetic field induces a current in the loop.
   - The setup shows how the rate of change of the magnetic field (\textbackslash{}( \textbackslash{}frac\{dB\}\{dt\} \textbackslash{})) affects the induced current (I) in the loop, as shown by the meter.

The diagram effectively illustrates the principles of electromagnetic induction and how current is generated in a loop within a magnetic field.

\#\# Dataset metadata

Electromagnetic Induction; Physical Model Grounding Reasoning, Multi-Formula Reasoning

\paragraph{Recorded answer/context.}
C: C: 0.048mA

\paragraph{Figure/image path.}
/root/proposal\_for\_physic/science-mango/phyx\_mini\_run/phyx\_data/test\_image/917.png

\paragraph{Formalization target.}
create a compiling Lean file with sorry bodies at `PhyXMiniProblems/problem\_phyx\_mini\_0917.lean`. The Lean declarations must preserve the physical quantities, dimensions or dimensional roles, figure labels, governing-law hypotheses, and final relation expressed by this problem.
Use Mathlib/Physlib names found through LeanExplore where available. If a physics API is missing, introduce faithful local abstractions rather than scalar placeholder aliases.

\begin{theorem}[Physics formalization target]
\label{thm:physics:phyx_mini_0917:target}
\lean{PhyXMiniProblems.ProblemPhyXMini0917.inducedCurrent_is_zeroPointZeroFourEightMilliamperes}
\uses{def:physics:phyx-mini-0917:phyxminiproblems-problemphyxmini0917-electriccurrentinmilliamperes, def:physics:phyx-mini-0917:phyxminiproblems-problemphyxmini0917-inductionloopsetup, def:physics:phyx-mini-0917:phyxminiproblems-problemphyxmini0917-matchesprimaryinductionfigure, def:physics:phyx-mini-0917:phyxminiproblems-problemphyxmini0917-matchesinductionproblemdescription, def:physics:phyx-mini-0917:phyxminiproblems-problemphyxmini0917-hasphysicalinductionparameters, def:physics:phyx-mini-0917:phyxminiproblems-problemphyxmini0917-satisfieselectromagneticinductionlaws, def:physics:phyx-mini-0917:phyxminiproblems-problemphyxmini0917-recordeddatasetanswer, def:physics:phyx-mini-0917:phyxminiproblems-problemphyxmini0917-answermatchesinducedcurrent, def:physics:phyx-mini-0917:phyxminiproblems-problemphyxmini0917-inducedcurrentrunsalongdrawnarrow}
The assigned autoformalize agent should translate this physics problem into Lean declarations in the covered file, with theorem and lemma proof bodies written as `by sorry`.
\end{theorem}
\begin{proof}
This is an autoformalization task, not a proof task. Produce statements that can later be proved by the physics prover without weakening the physical model.
\end{proof}

% --- Archon declaration topology begin ---
\section{Declaration topology}
\begin{definition}[area Dimension]
  \label{def:physics:phyx-mini-0917:phyxminiproblems-problemphyxmini0917-areadimension}
  \lean{PhyXMiniProblems.ProblemPhyXMini0917.areaDimension}
  Physical dimension L² of loop area
\end{definition}
\begin{proof}
  This is the declared model definition of area Dimension.
\end{proof}
\begin{definition}[electrical Resistance Dimension]
  \label{def:physics:phyx-mini-0917:phyxminiproblems-problemphyxmini0917-electricalresistancedimension}
  \lean{PhyXMiniProblems.ProblemPhyXMini0917.electricalResistanceDimension}
  Physical dimension M L² T⁻¹ C⁻² of electrical resistance
\end{definition}
\begin{proof}
  This is the declared model definition of electrical Resistance Dimension.
\end{proof}
\begin{definition}[magnetic Flux Density Dimension]
  \label{def:physics:phyx-mini-0917:phyxminiproblems-problemphyxmini0917-magneticfluxdensitydimension}
  \lean{PhyXMiniProblems.ProblemPhyXMini0917.magneticFluxDensityDimension}
  Physical dimension M T⁻¹ C⁻¹ of magnetic flux density
\end{definition}
\begin{proof}
  This is the declared model definition of magnetic Flux Density Dimension.
\end{proof}
\begin{definition}[magnetic Flux Density Rate Dimension]
  \label{def:physics:phyx-mini-0917:phyxminiproblems-problemphyxmini0917-magneticfluxdensityratedimension}
  \lean{PhyXMiniProblems.ProblemPhyXMini0917.magneticFluxDensityRateDimension}
  Physical dimension M T⁻² C⁻¹ of a magnetic-flux-density rate
\end{definition}
\begin{proof}
  This is the declared model definition of magnetic Flux Density Rate Dimension.
\end{proof}
\begin{definition}[magnetic Flux Rate Dimension]
  \label{def:physics:phyx-mini-0917:phyxminiproblems-problemphyxmini0917-magneticfluxratedimension}
  \lean{PhyXMiniProblems.ProblemPhyXMini0917.magneticFluxRateDimension}
  Physical dimension M L² T⁻² C⁻¹ of magnetic-flux rate
\end{definition}
\begin{proof}
  This is the declared model definition of magnetic Flux Rate Dimension.
\end{proof}
\begin{definition}[electromotive Force Dimension]
  \label{def:physics:phyx-mini-0917:phyxminiproblems-problemphyxmini0917-electromotiveforcedimension}
  \lean{PhyXMiniProblems.ProblemPhyXMini0917.electromotiveForceDimension}
  Physical dimension M L² T⁻² C⁻¹ of electromotive force
\end{definition}
\begin{proof}
  This is the declared model definition of electromotive Force Dimension.
\end{proof}
\begin{definition}[electric Current Dimension]
  \label{def:physics:phyx-mini-0917:phyxminiproblems-problemphyxmini0917-electriccurrentdimension}
  \lean{PhyXMiniProblems.ProblemPhyXMini0917.electricCurrentDimension}
  Physical dimension C T⁻¹ of electric current
\end{definition}
\begin{proof}
  This is the declared model definition of electric Current Dimension.
\end{proof}
\begin{definition}[Area Magnitude]
  \label{def:physics:phyx-mini-0917:phyxminiproblems-problemphyxmini0917-areamagnitude}
  \lean{PhyXMiniProblems.ProblemPhyXMini0917.AreaMagnitude}
  \uses{def:physics:phyx-mini-0917:phyxminiproblems-problemphyxmini0917-areadimension}
  A nonnegative, unit-independent physical area
\end{definition}
\begin{proof}
  This is the declared model definition of Area Magnitude.
\end{proof}
\begin{definition}[Electrical Resistance Magnitude]
  \label{def:physics:phyx-mini-0917:phyxminiproblems-problemphyxmini0917-electricalresistancemagnitude}
  \lean{PhyXMiniProblems.ProblemPhyXMini0917.ElectricalResistanceMagnitude}
  \uses{def:physics:phyx-mini-0917:phyxminiproblems-problemphyxmini0917-electricalresistancedimension}
  A nonnegative, unit-independent electrical resistance
\end{definition}
\begin{proof}
  This is the declared model definition of Electrical Resistance Magnitude.
\end{proof}
\begin{definition}[Magnetic Flux Density Magnitude]
  \label{def:physics:phyx-mini-0917:phyxminiproblems-problemphyxmini0917-magneticfluxdensitymagnitude}
  \lean{PhyXMiniProblems.ProblemPhyXMini0917.MagneticFluxDensityMagnitude}
  \uses{def:physics:phyx-mini-0917:phyxminiproblems-problemphyxmini0917-magneticfluxdensitydimension}
  A nonnegative, unit-independent magnetic-flux-density magnitude
\end{definition}
\begin{proof}
  This is the declared model definition of Magnetic Flux Density Magnitude.
\end{proof}
\begin{definition}[Magnetic Flux Density Rate Magnitude]
  \label{def:physics:phyx-mini-0917:phyxminiproblems-problemphyxmini0917-magneticfluxdensityratemagnitude}
  \lean{PhyXMiniProblems.ProblemPhyXMini0917.MagneticFluxDensityRateMagnitude}
  \uses{def:physics:phyx-mini-0917:phyxminiproblems-problemphyxmini0917-magneticfluxdensityratedimension}
  A nonnegative, unit-independent magnetic-flux-density rate
\end{definition}
\begin{proof}
  This is the declared model definition of Magnetic Flux Density Rate Magnitude.
\end{proof}
\begin{definition}[Signed Magnetic Flux Rate]
  \label{def:physics:phyx-mini-0917:phyxminiproblems-problemphyxmini0917-signedmagneticfluxrate}
  \lean{PhyXMiniProblems.ProblemPhyXMini0917.SignedMagneticFluxRate}
  \uses{def:physics:phyx-mini-0917:phyxminiproblems-problemphyxmini0917-magneticfluxratedimension}
  A signed, unit-independent magnetic-flux rate
\end{definition}
\begin{proof}
  This is the declared model definition of Signed Magnetic Flux Rate.
\end{proof}
\begin{definition}[Electromotive Force]
  \label{def:physics:phyx-mini-0917:phyxminiproblems-problemphyxmini0917-electromotiveforce}
  \lean{PhyXMiniProblems.ProblemPhyXMini0917.ElectromotiveForce}
  \uses{def:physics:phyx-mini-0917:phyxminiproblems-problemphyxmini0917-electromotiveforcedimension}
  A signed, unit-independent electromotive force
\end{definition}
\begin{proof}
  This is the declared model definition of Electromotive Force.
\end{proof}
\begin{definition}[Electric Current]
  \label{def:physics:phyx-mini-0917:phyxminiproblems-problemphyxmini0917-electriccurrent}
  \lean{PhyXMiniProblems.ProblemPhyXMini0917.ElectricCurrent}
  \uses{def:physics:phyx-mini-0917:phyxminiproblems-problemphyxmini0917-electriccurrentdimension}
  A signed, unit-independent electric current
\end{definition}
\begin{proof}
  This is the declared model definition of Electric Current.
\end{proof}
\begin{definition}[area Readout]
  \label{def:physics:phyx-mini-0917:phyxminiproblems-problemphyxmini0917-areareadout}
  \lean{PhyXMiniProblems.ProblemPhyXMini0917.areaReadout}
  \uses{def:physics:phyx-mini-0917:phyxminiproblems-problemphyxmini0917-areamagnitude}
  Read a physical area in the coherent area unit selected by units
\end{definition}
\begin{proof}
  This is the declared model definition of area Readout.
\end{proof}
\begin{definition}[resistance Readout]
  \label{def:physics:phyx-mini-0917:phyxminiproblems-problemphyxmini0917-resistancereadout}
  \lean{PhyXMiniProblems.ProblemPhyXMini0917.resistanceReadout}
  \uses{def:physics:phyx-mini-0917:phyxminiproblems-problemphyxmini0917-electricalresistancemagnitude}
  Read resistance in the coherent resistance unit selected by units
\end{definition}
\begin{proof}
  This is the declared model definition of resistance Readout.
\end{proof}
\begin{definition}[magnetic Flux Density Readout]
  \label{def:physics:phyx-mini-0917:phyxminiproblems-problemphyxmini0917-magneticfluxdensityreadout}
  \lean{PhyXMiniProblems.ProblemPhyXMini0917.magneticFluxDensityReadout}
  \uses{def:physics:phyx-mini-0917:phyxminiproblems-problemphyxmini0917-magneticfluxdensitymagnitude}
  Read magnetic flux density in the coherent unit selected by units
\end{definition}
\begin{proof}
  This is the declared model definition of magnetic Flux Density Readout.
\end{proof}
\begin{definition}[magnetic Flux Density Rate Readout]
  \label{def:physics:phyx-mini-0917:phyxminiproblems-problemphyxmini0917-magneticfluxdensityratereadout}
  \lean{PhyXMiniProblems.ProblemPhyXMini0917.magneticFluxDensityRateReadout}
  \uses{def:physics:phyx-mini-0917:phyxminiproblems-problemphyxmini0917-magneticfluxdensityratemagnitude}
  Read magnetic-flux-density rate in the coherent unit selected by units
\end{definition}
\begin{proof}
  This is the declared model definition of magnetic Flux Density Rate Readout.
\end{proof}
\begin{definition}[signed Magnetic Flux Rate Readout]
  \label{def:physics:phyx-mini-0917:phyxminiproblems-problemphyxmini0917-signedmagneticfluxratereadout}
  \lean{PhyXMiniProblems.ProblemPhyXMini0917.signedMagneticFluxRateReadout}
  \uses{def:physics:phyx-mini-0917:phyxminiproblems-problemphyxmini0917-signedmagneticfluxrate}
  Read signed magnetic-flux rate in the coherent unit selected by units
\end{definition}
\begin{proof}
  This is the declared model definition of signed Magnetic Flux Rate Readout.
\end{proof}
\begin{definition}[electromotive Force Readout]
  \label{def:physics:phyx-mini-0917:phyxminiproblems-problemphyxmini0917-electromotiveforcereadout}
  \lean{PhyXMiniProblems.ProblemPhyXMini0917.electromotiveForceReadout}
  \uses{def:physics:phyx-mini-0917:phyxminiproblems-problemphyxmini0917-electromotiveforce}
  Read electromotive force in the coherent unit selected by units
\end{definition}
\begin{proof}
  This is the declared model definition of electromotive Force Readout.
\end{proof}
\begin{definition}[electric Current Readout]
  \label{def:physics:phyx-mini-0917:phyxminiproblems-problemphyxmini0917-electriccurrentreadout}
  \lean{PhyXMiniProblems.ProblemPhyXMini0917.electricCurrentReadout}
  \uses{def:physics:phyx-mini-0917:phyxminiproblems-problemphyxmini0917-electriccurrent}
  Read signed electric current in the coherent unit selected by units
\end{definition}
\begin{proof}
  This is the declared model definition of electric Current Readout.
\end{proof}
\begin{definition}[area In Square Meters]
  \label{def:physics:phyx-mini-0917:phyxminiproblems-problemphyxmini0917-areainsquaremeters}
  \lean{PhyXMiniProblems.ProblemPhyXMini0917.areaInSquareMeters}
  \uses{def:physics:phyx-mini-0917:phyxminiproblems-problemphyxmini0917-areamagnitude, def:physics:phyx-mini-0917:phyxminiproblems-problemphyxmini0917-areareadout}
  Coherent-SI square-metre readout of loop area
\end{definition}
\begin{proof}
  This is the declared model definition of area In Square Meters.
\end{proof}
\begin{definition}[area In Square Centimeters]
  \label{def:physics:phyx-mini-0917:phyxminiproblems-problemphyxmini0917-areainsquarecentimeters}
  \lean{PhyXMiniProblems.ProblemPhyXMini0917.areaInSquareCentimeters}
  \uses{def:physics:phyx-mini-0917:phyxminiproblems-problemphyxmini0917-areamagnitude, def:physics:phyx-mini-0917:phyxminiproblems-problemphyxmini0917-areainsquaremeters}
  Square-centimetre readout used in the printed area label
\end{definition}
\begin{proof}
  This is the declared model definition of area In Square Centimeters.
\end{proof}
\begin{definition}[resistance In Ohms]
  \label{def:physics:phyx-mini-0917:phyxminiproblems-problemphyxmini0917-resistanceinohms}
  \lean{PhyXMiniProblems.ProblemPhyXMini0917.resistanceInOhms}
  \uses{def:physics:phyx-mini-0917:phyxminiproblems-problemphyxmini0917-electricalresistancemagnitude, def:physics:phyx-mini-0917:phyxminiproblems-problemphyxmini0917-resistancereadout}
  Coherent-SI ohm readout of resistance
\end{definition}
\begin{proof}
  This is the declared model definition of resistance In Ohms.
\end{proof}
\begin{definition}[magnetic Flux Density In Teslas]
  \label{def:physics:phyx-mini-0917:phyxminiproblems-problemphyxmini0917-magneticfluxdensityinteslas}
  \lean{PhyXMiniProblems.ProblemPhyXMini0917.magneticFluxDensityInTeslas}
  \uses{def:physics:phyx-mini-0917:phyxminiproblems-problemphyxmini0917-magneticfluxdensitymagnitude, def:physics:phyx-mini-0917:phyxminiproblems-problemphyxmini0917-magneticfluxdensityreadout}
  Coherent-SI tesla readout of magnetic-flux-density magnitude
\end{definition}
\begin{proof}
  This is the declared model definition of magnetic Flux Density In Teslas.
\end{proof}
\begin{definition}[magnetic Flux Density Rate In Teslas Per Second]
  \label{def:physics:phyx-mini-0917:phyxminiproblems-problemphyxmini0917-magneticfluxdensityrateinteslaspersecond}
  \lean{PhyXMiniProblems.ProblemPhyXMini0917.magneticFluxDensityRateInTeslasPerSecond}
  \uses{def:physics:phyx-mini-0917:phyxminiproblems-problemphyxmini0917-magneticfluxdensityratemagnitude, def:physics:phyx-mini-0917:phyxminiproblems-problemphyxmini0917-magneticfluxdensityratereadout}
  Coherent-SI tesla-per-second readout of field-magnitude rate
\end{definition}
\begin{proof}
  This is the declared model definition of magnetic Flux Density Rate In Teslas Per Second.
\end{proof}
\begin{definition}[signed Magnetic Flux Rate In Webers Per Second]
  \label{def:physics:phyx-mini-0917:phyxminiproblems-problemphyxmini0917-signedmagneticfluxrateinweberspersecond}
  \lean{PhyXMiniProblems.ProblemPhyXMini0917.signedMagneticFluxRateInWebersPerSecond}
  \uses{def:physics:phyx-mini-0917:phyxminiproblems-problemphyxmini0917-signedmagneticfluxrate, def:physics:phyx-mini-0917:phyxminiproblems-problemphyxmini0917-signedmagneticfluxratereadout}
  Coherent-SI weber-per-second readout of signed magnetic-flux rate
\end{definition}
\begin{proof}
  This is the declared model definition of signed Magnetic Flux Rate In Webers Per Second.
\end{proof}
\begin{definition}[electromotive Force In Volts]
  \label{def:physics:phyx-mini-0917:phyxminiproblems-problemphyxmini0917-electromotiveforceinvolts}
  \lean{PhyXMiniProblems.ProblemPhyXMini0917.electromotiveForceInVolts}
  \uses{def:physics:phyx-mini-0917:phyxminiproblems-problemphyxmini0917-electromotiveforce, def:physics:phyx-mini-0917:phyxminiproblems-problemphyxmini0917-electromotiveforcereadout}
  Coherent-SI volt readout of electromotive force
\end{definition}
\begin{proof}
  This is the declared model definition of electromotive Force In Volts.
\end{proof}
\begin{definition}[electric Current In Amperes]
  \label{def:physics:phyx-mini-0917:phyxminiproblems-problemphyxmini0917-electriccurrentinamperes}
  \lean{PhyXMiniProblems.ProblemPhyXMini0917.electricCurrentInAmperes}
  \uses{def:physics:phyx-mini-0917:phyxminiproblems-problemphyxmini0917-electriccurrent, def:physics:phyx-mini-0917:phyxminiproblems-problemphyxmini0917-electriccurrentreadout}
  Coherent-SI ampere readout, signed relative to the drawn arrow
\end{definition}
\begin{proof}
  This is the declared model definition of electric Current In Amperes.
\end{proof}
\begin{definition}[electric Current In Milliamperes]
  \label{def:physics:phyx-mini-0917:phyxminiproblems-problemphyxmini0917-electriccurrentinmilliamperes}
  \lean{PhyXMiniProblems.ProblemPhyXMini0917.electricCurrentInMilliamperes}
  \uses{def:physics:phyx-mini-0917:phyxminiproblems-problemphyxmini0917-electriccurrent, def:physics:phyx-mini-0917:phyxminiproblems-problemphyxmini0917-electriccurrentinamperes}
  Milliampere readout, signed relative to the drawn arrow
\end{definition}
\begin{proof}
  This is the declared model definition of electric Current In Milliamperes.
\end{proof}
\begin{definition}[Magnetic Pole]
  \label{def:physics:phyx-mini-0917:phyxminiproblems-problemphyxmini0917-magneticpole}
  \lean{PhyXMiniProblems.ProblemPhyXMini0917.MagneticPole}
  Labels printed on the two electromagnet pole pieces
\end{definition}
\begin{proof}
  This is the declared model definition of Magnetic Pole.
\end{proof}
\begin{definition}[Pole Position]
  \label{def:physics:phyx-mini-0917:phyxminiproblems-problemphyxmini0917-poleposition}
  \lean{PhyXMiniProblems.ProblemPhyXMini0917.PolePosition}
  Vertical locations of the two pole pieces in the primary raster
\end{definition}
\begin{proof}
  This is the declared model definition of Pole Position.
\end{proof}
\begin{definition}[Loop Point]
  \label{def:physics:phyx-mini-0917:phyxminiproblems-problemphyxmini0917-looppoint}
  \lean{PhyXMiniProblems.ProblemPhyXMini0917.LoopPoint}
  The two lettered points where the external circuit meets the loop
\end{definition}
\begin{proof}
  This is the declared model definition of Loop Point.
\end{proof}
\begin{definition}[Vertical Direction]
  \label{def:physics:phyx-mini-0917:phyxminiproblems-problemphyxmini0917-verticaldirection}
  \lean{PhyXMiniProblems.ProblemPhyXMini0917.VerticalDirection}
  Vertical directions available to the field and area-vector arrows
\end{definition}
\begin{proof}
  This is the declared model definition of Vertical Direction.
\end{proof}
\begin{definition}[Loop Traversal Direction]
  \label{def:physics:phyx-mini-0917:phyxminiproblems-problemphyxmini0917-looptraversaldirection}
  \lean{PhyXMiniProblems.ProblemPhyXMini0917.LoopTraversalDirection}
  Traversal directions at the visible front part of the loop
\end{definition}
\begin{proof}
  This is the declared model definition of Loop Traversal Direction.
\end{proof}
\begin{definition}[Axial Alignment]
  \label{def:physics:phyx-mini-0917:phyxminiproblems-problemphyxmini0917-axialalignment}
  \lean{PhyXMiniProblems.ProblemPhyXMini0917.AxialAlignment}
  Parallel or antiparallel alignment of two axial directions
\end{definition}
\begin{proof}
  This is the declared model definition of Axial Alignment.
\end{proof}
\begin{definition}[Loop Topology]
  \label{def:physics:phyx-mini-0917:phyxminiproblems-problemphyxmini0917-looptopology}
  \lean{PhyXMiniProblems.ProblemPhyXMini0917.LoopTopology}
  Topology and conductor count of the loop shown in the problem
\end{definition}
\begin{proof}
  This is the declared model definition of Loop Topology.
\end{proof}
\begin{definition}[Loop Motion Model]
  \label{def:physics:phyx-mini-0917:phyxminiproblems-problemphyxmini0917-loopmotionmodel}
  \lean{PhyXMiniProblems.ProblemPhyXMini0917.LoopMotionModel}
  Mechanical state relevant to the motional-emf term
\end{definition}
\begin{proof}
  This is the declared model definition of Loop Motion Model.
\end{proof}
\begin{definition}[Field Spatial Uniformity]
  \label{def:physics:phyx-mini-0917:phyxminiproblems-problemphyxmini0917-fieldspatialuniformity}
  \lean{PhyXMiniProblems.ProblemPhyXMini0917.FieldSpatialUniformity}
  Spatial idealization stated for the field at each instant
\end{definition}
\begin{proof}
  This is the declared model definition of Field Spatial Uniformity.
\end{proof}
\begin{definition}[Meter Kind]
  \label{def:physics:phyx-mini-0917:phyxminiproblems-problemphyxmini0917-meterkind}
  \lean{PhyXMiniProblems.ProblemPhyXMini0917.MeterKind}
  Kind of meter connected into the single series circuit
\end{definition}
\begin{proof}
  This is the declared model definition of Meter Kind.
\end{proof}
\begin{definition}[expected Pole At]
  \label{def:physics:phyx-mini-0917:phyxminiproblems-problemphyxmini0917-expectedpoleat}
  \lean{PhyXMiniProblems.ProblemPhyXMini0917.expectedPoleAt}
  \uses{def:physics:phyx-mini-0917:phyxminiproblems-problemphyxmini0917-magneticpole, def:physics:phyx-mini-0917:phyxminiproblems-problemphyxmini0917-poleposition}
  Pole label expected at each vertical location in image 917.png
\end{definition}
\begin{proof}
  This is the declared model definition of expected Pole At.
\end{proof}
\begin{definition}[right Hand Normal]
  \label{def:physics:phyx-mini-0917:phyxminiproblems-problemphyxmini0917-righthandnormal}
  \lean{PhyXMiniProblems.ProblemPhyXMini0917.rightHandNormal}
  \uses{def:physics:phyx-mini-0917:phyxminiproblems-problemphyxmini0917-verticaldirection, def:physics:phyx-mini-0917:phyxminiproblems-problemphyxmini0917-looptraversaldirection}
  Right-hand normal generated by a chosen positive loop traversal
\end{definition}
\begin{proof}
  This is the declared model definition of right Hand Normal.
\end{proof}
\begin{definition}[axial Alignment]
  \label{def:physics:phyx-mini-0917:phyxminiproblems-problemphyxmini0917-axialalignment-value}
  \lean{PhyXMiniProblems.ProblemPhyXMini0917.axialAlignment}
  \uses{def:physics:phyx-mini-0917:phyxminiproblems-problemphyxmini0917-verticaldirection, def:physics:phyx-mini-0917:phyxminiproblems-problemphyxmini0917-axialalignment}
  Whether two vertical axial directions are parallel or antiparallel
\end{definition}
\begin{proof}
  This is the declared model definition of axial Alignment.
\end{proof}
\begin{definition}[cosine]
  \label{def:physics:phyx-mini-0917:phyxminiproblems-problemphyxmini0917-axialalignment-cosine}
  \lean{PhyXMiniProblems.ProblemPhyXMini0917.AxialAlignment.cosine}
  \uses{def:physics:phyx-mini-0917:phyxminiproblems-problemphyxmini0917-axialalignment}
  Cosine factor for the only two alignments appearing in this problem
\end{definition}
\begin{proof}
  This is the declared model definition of cosine.
\end{proof}
\begin{definition}[Electromagnetic Induction Figure]
  \label{def:physics:phyx-mini-0917:phyxminiproblems-problemphyxmini0917-electromagneticinductionfigure}
  \lean{PhyXMiniProblems.ProblemPhyXMini0917.ElectromagneticInductionFigure}
  \uses{def:physics:phyx-mini-0917:phyxminiproblems-problemphyxmini0917-magneticpole, def:physics:phyx-mini-0917:phyxminiproblems-problemphyxmini0917-poleposition, def:physics:phyx-mini-0917:phyxminiproblems-problemphyxmini0917-looppoint, def:physics:phyx-mini-0917:phyxminiproblems-problemphyxmini0917-verticaldirection, def:physics:phyx-mini-0917:phyxminiproblems-problemphyxmini0917-looptraversaldirection, def:physics:phyx-mini-0917:phyxminiproblems-problemphyxmini0917-meterkind}
  Literal geometric and numerical content transcribed from the primary raster. The purple arrow chooses the positive current sign; it does not assert which way the induced current actually flows
\end{definition}
\begin{proof}
  This is the declared model definition of Electromagnetic Induction Figure.
\end{proof}
\begin{definition}[Induction Loop Setup]
  \label{def:physics:phyx-mini-0917:phyxminiproblems-problemphyxmini0917-inductionloopsetup}
  \lean{PhyXMiniProblems.ProblemPhyXMini0917.InductionLoopSetup}
  \uses{def:physics:phyx-mini-0917:phyxminiproblems-problemphyxmini0917-areamagnitude, def:physics:phyx-mini-0917:phyxminiproblems-problemphyxmini0917-electricalresistancemagnitude, def:physics:phyx-mini-0917:phyxminiproblems-problemphyxmini0917-magneticfluxdensitymagnitude, def:physics:phyx-mini-0917:phyxminiproblems-problemphyxmini0917-magneticfluxdensityratemagnitude, def:physics:phyx-mini-0917:phyxminiproblems-problemphyxmini0917-signedmagneticfluxrate, def:physics:phyx-mini-0917:phyxminiproblems-problemphyxmini0917-electromotiveforce, def:physics:phyx-mini-0917:phyxminiproblems-problemphyxmini0917-electriccurrent, def:physics:phyx-mini-0917:phyxminiproblems-problemphyxmini0917-verticaldirection, def:physics:phyx-mini-0917:phyxminiproblems-problemphyxmini0917-looptraversaldirection, def:physics:phyx-mini-0917:phyxminiproblems-problemphyxmini0917-looptopology, def:physics:phyx-mini-0917:phyxminiproblems-problemphyxmini0917-loopmotionmodel, def:physics:phyx-mini-0917:phyxminiproblems-problemphyxmini0917-fieldspatialuniformity, def:physics:phyx-mini-0917:phyxminiproblems-problemphyxmini0917-electromagneticinductionfigure}
  Independent physical objects and observables. In particular, induced emf and current are not defined from any answer choice or numerical target
\end{definition}
\begin{proof}
  This is the declared model definition of Induction Loop Setup.
\end{proof}
\begin{definition}[Matches Primary Induction Figure]
  \label{def:physics:phyx-mini-0917:phyxminiproblems-problemphyxmini0917-matchesprimaryinductionfigure}
  \lean{PhyXMiniProblems.ProblemPhyXMini0917.MatchesPrimaryInductionFigure}
  \uses{def:physics:phyx-mini-0917:phyxminiproblems-problemphyxmini0917-expectedpoleat, def:physics:phyx-mini-0917:phyxminiproblems-problemphyxmini0917-inductionloopsetup}
  Exact labels, arrows, components, and printed values in image 917.png
\end{definition}
\begin{proof}
  This is the declared model definition of Matches Primary Induction Figure.
\end{proof}
\begin{definition}[Matches Induction Problem Description]
  \label{def:physics:phyx-mini-0917:phyxminiproblems-problemphyxmini0917-matchesinductionproblemdescription}
  \lean{PhyXMiniProblems.ProblemPhyXMini0917.MatchesInductionProblemDescription}
  \uses{def:physics:phyx-mini-0917:phyxminiproblems-problemphyxmini0917-areainsquaremeters, def:physics:phyx-mini-0917:phyxminiproblems-problemphyxmini0917-areainsquarecentimeters, def:physics:phyx-mini-0917:phyxminiproblems-problemphyxmini0917-resistanceinohms, def:physics:phyx-mini-0917:phyxminiproblems-problemphyxmini0917-magneticfluxdensityrateinteslaspersecond, def:physics:phyx-mini-0917:phyxminiproblems-problemphyxmini0917-inductionloopsetup}
  Qualitative setup and calibration of physical quantities to raster data
\end{definition}
\begin{proof}
  This is the declared model definition of Matches Induction Problem Description.
\end{proof}
\begin{definition}[Has Physical Induction Parameters]
  \label{def:physics:phyx-mini-0917:phyxminiproblems-problemphyxmini0917-hasphysicalinductionparameters}
  \lean{PhyXMiniProblems.ProblemPhyXMini0917.HasPhysicalInductionParameters}
  \uses{def:physics:phyx-mini-0917:phyxminiproblems-problemphyxmini0917-areainsquaremeters, def:physics:phyx-mini-0917:phyxminiproblems-problemphyxmini0917-resistanceinohms, def:physics:phyx-mini-0917:phyxminiproblems-problemphyxmini0917-magneticfluxdensityrateinteslaspersecond, def:physics:phyx-mini-0917:phyxminiproblems-problemphyxmini0917-inductionloopsetup}
  Positivity conditions selecting the ordinary passive-circuit branch
\end{definition}
\begin{proof}
  This is the declared model definition of Has Physical Induction Parameters.
\end{proof}
\begin{definition}[Satisfies Electromagnetic Induction Laws]
  \label{def:physics:phyx-mini-0917:phyxminiproblems-problemphyxmini0917-satisfieselectromagneticinductionlaws}
  \lean{PhyXMiniProblems.ProblemPhyXMini0917.SatisfiesElectromagneticInductionLaws}
  \uses{def:physics:phyx-mini-0917:phyxminiproblems-problemphyxmini0917-areainsquaremeters, def:physics:phyx-mini-0917:phyxminiproblems-problemphyxmini0917-resistanceinohms, def:physics:phyx-mini-0917:phyxminiproblems-problemphyxmini0917-magneticfluxdensityinteslas, def:physics:phyx-mini-0917:phyxminiproblems-problemphyxmini0917-magneticfluxdensityrateinteslaspersecond, def:physics:phyx-mini-0917:phyxminiproblems-problemphyxmini0917-signedmagneticfluxrateinweberspersecond, def:physics:phyx-mini-0917:phyxminiproblems-problemphyxmini0917-electromotiveforceinvolts, def:physics:phyx-mini-0917:phyxminiproblems-problemphyxmini0917-electriccurrentinamperes, def:physics:phyx-mini-0917:phyxminiproblems-problemphyxmini0917-righthandnormal, def:physics:phyx-mini-0917:phyxminiproblems-problemphyxmini0917-axialalignment-value, def:physics:phyx-mini-0917:phyxminiproblems-problemphyxmini0917-inductionloopsetup}
  Faraday induction and the series-circuit constitutive law. The flux-rate law uses the normal determined by the current sign reference, which is distinct from the independently displayed upward area vector
\end{definition}
\begin{proof}
  This is the declared model definition of Satisfies Electromagnetic Induction Laws.
\end{proof}
\begin{definition}[Answer Choice]
  \label{def:physics:phyx-mini-0917:phyxminiproblems-problemphyxmini0917-answerchoice}
  \lean{PhyXMiniProblems.ProblemPhyXMini0917.AnswerChoice}
  The four signed milliampere answers printed with the problem
\end{definition}
\begin{proof}
  This is the declared model definition of Answer Choice.
\end{proof}
\begin{definition}[displayed Current In Milliamperes]
  \label{def:physics:phyx-mini-0917:phyxminiproblems-problemphyxmini0917-answerchoice-displayedcurrentinmilliamperes}
  \lean{PhyXMiniProblems.ProblemPhyXMini0917.AnswerChoice.displayedCurrentInMilliamperes}
  \uses{def:physics:phyx-mini-0917:phyxminiproblems-problemphyxmini0917-answerchoice}
  Signed current in milliamperes printed beside an answer choice
\end{definition}
\begin{proof}
  This is the declared model definition of displayed Current In Milliamperes.
\end{proof}
\begin{definition}[recorded Dataset Answer]
  \label{def:physics:phyx-mini-0917:phyxminiproblems-problemphyxmini0917-recordeddatasetanswer}
  \lean{PhyXMiniProblems.ProblemPhyXMini0917.recordedDatasetAnswer}
  \uses{def:physics:phyx-mini-0917:phyxminiproblems-problemphyxmini0917-answerchoice}
  Answer label recorded in the supplied dataset
\end{definition}
\begin{proof}
  This is the declared model definition of recorded Dataset Answer.
\end{proof}
\begin{definition}[Answer Matches Induced Current]
  \label{def:physics:phyx-mini-0917:phyxminiproblems-problemphyxmini0917-answermatchesinducedcurrent}
  \lean{PhyXMiniProblems.ProblemPhyXMini0917.AnswerMatchesInducedCurrent}
  \uses{def:physics:phyx-mini-0917:phyxminiproblems-problemphyxmini0917-electriccurrentinmilliamperes, def:physics:phyx-mini-0917:phyxminiproblems-problemphyxmini0917-inductionloopsetup, def:physics:phyx-mini-0917:phyxminiproblems-problemphyxmini0917-answerchoice}
  A choice matches the independently modeled induced current
\end{definition}
\begin{proof}
  This is the declared model definition of Answer Matches Induced Current.
\end{proof}
\begin{definition}[Induced Current Runs Along Drawn Arrow]
  \label{def:physics:phyx-mini-0917:phyxminiproblems-problemphyxmini0917-inducedcurrentrunsalongdrawnarrow}
  \lean{PhyXMiniProblems.ProblemPhyXMini0917.InducedCurrentRunsAlongDrawnArrow}
  \uses{def:physics:phyx-mini-0917:phyxminiproblems-problemphyxmini0917-electriccurrentinamperes, def:physics:phyx-mini-0917:phyxminiproblems-problemphyxmini0917-inductionloopsetup}
  Positive signed current means motion along the purple reference arrow
\end{definition}
\begin{proof}
  This is the declared model definition of Induced Current Runs Along Drawn Arrow.
\end{proof}
% --- Archon declaration topology end ---
% --- Archon physics formalization source end ---
